\documentclass[10pt,a4paper]{amsart}
\usepackage[margin=19mm]{geometry}
\usepackage{amsmath,amssymb,mathtools}
\usepackage[scr=rsfs,cal=euler]{mathalfa}
\usepackage{xfrac}
\usepackage[unicode,hidelinks,bookmarks=false]{hyperref}
\newcommand{\ie}{i.e.\ }
\newcommand{\cf}{cf.\ }
\newcommand{\resp}{respectively\ }
\newcommand{\Subsection}{\subsection}
\providecommand{\nobreakdash}{\nobreak\hbox{-}}
\setlength{\emergencystretch}{2em}
\multlinegap0pt
\usepackage{xcolor}
\usepackage{amssymb}
\usepackage{bm}
\usepackage[normalem]{ulem}
\usepackage{mathtools}
\usepackage{tikz}
\newtheorem{prop}{Proposition}
\def\O{\mathcal{O}}
\def\m{\mu}
\def\mod{\operatorname{mod}}
\def\wAm{\widehat{A}(m)}
\def\zm{\mathbb{Z}/m\mathbb{Z}}
\title{Spectral correspondence for cyclic Higgs bundles}
\author{JIA CHOON LEE and Ana Pe\'on-Nieto}
\date{Source: arXiv:2605.03249v2 (CC BY 4.0)}
\begin{document}
\maketitle

\noindent\emph{Source and licence.} Spectral correspondence for cyclic Higgs bundles. Version arXiv:2605.03249v2 (CC BY 4.0), distributed by arXiv under the Creative Commons Attribution 4.0 International licence, http://creativecommons.org/licenses/by/4.0/, as recorded in the arXiv licence metadata for this exact version and shown on the abstract page https://arxiv.org/abs/2605.03249v2. Full licence notice: \texttt{licenses/arxiv-2605.03249-CC-BY-4.0.txt}.\par\smallskip

\noindent\textit{Editorial scope.} Counted unit: the article's prop prop:equiv3 (source0.tex lines 523--527). The article's complete original proof of it is delivered with it (source0.tex lines 529--547, one \texttt{proof} environment of 19 lines). No mathematics is written, repaired or re-derived: the statement and the proof are the article's own lines, in the article's own notation, and no retained line is substituted or re-flowed.  Retained uncounted as the source-local context the proof needs: lines 375--377.  Nothing was omitted from any retained span, so the package carries no omission record.  No retained line contains a citation, so the wrapper carries no bibliography and no bibliography entry.  No retained line contains a figure environment, an image-inclusion command or drawing code, so no image is delivered and the package declares no asset dependency.  Editorial formatting only, in addition: an AMS \texttt{amsart} preamble that restates the article's own declarations and macros used by the retained lines, namely the environments \texttt{prop} on separate counters, together with the standard packages those lines need.  The retained environments print in this document as Proposition 1 in order of appearance, whereas the article prints the same items with the numbering of its own full document.  The complete native source of the article as distributed in the arXiv e-print for this exact version is bundled unchanged as \texttt{sources/199779/source0.tex}.
\begin{prop}\label{prop:equiv1}
There is an equivalence between the category of $M$-twisted $Q$-quiver bundles on $C$ and the category of coherent right $A$-modules $\mod_C\!\text{-}A$.
\end{prop}

\begin{prop}\label{prop:equiv3}
There is an equivalence between the category $\mod_Y\!\text{-}\wAm$ of coherent right $\wAm$-modules on $Y$ and the category of $M_Y$-twisted $Q(m)$-quiver sheaves $(\{F_i\}_{i\in \zm}, \{\psi_i\}_{i\in \zm} )$ with the additional condition that for each $i\in \zm$, the composition of arrows around the loop (suppressing the twist by the identity on $\pi^*K_C^{-1}$) yields:
\[\psi_{i-1}\circ \cdots \circ \psi_i = \operatorname{Id}_{F_i}\otimes_{\O_Y} \lambda  \in Hom_{\O_Y}(F_i\otimes_{\O_Y} \pi^*K_C^{-m}, F_i )\]
where the subcript of $\psi_i$ is understood modulo $m$.
\end{prop}

\begin{proof}
    The category $\mod_Y\!\text{-}\wAm$ is a full subcategory of $\mod_Y\!\text{-}A_Y(m)$ consisting of $A_Y(m)$-modules that annihilates the ideal $I$. By Proposition \ref{prop:equiv1}, an object $\m F\in \mod_Y\!\text{-}A_Y(m)$ corresponds to a $M_Y$-twisted $Q(m)$-quiver sheaf $(\{F_i\}_{i\in \zm}, \{\psi_i\}_{i\in \zm})$. It remains to determine what condition is imposed on this quiver sheaf by requiring it to annihilate $I$. 

    Recall that $I=Im(\nu_1-\nu_2)$. 
    The map $\nu_1: \pi^*K_C^{-m}\to A_Y(m)$ corresponding to loops in the path algebra induces
    \begin{align*}
    \m  F\otimes_{\O_Y} \pi^*K_C^{-m} &\to \m F\otimes_{\O_Y} A_Y(m) \to \m F
     %f\otimes z &\mapsto  f \cdot \left(\sum_{i\in \zm}\iota_i(z)\right)\mapsto  f \cdot \left(\sum_{i\in \zm}\iota_i(z)\right)
    \end{align*}
    By restricting this map to the component $F_i$, we get
    \[\psi_{i-1}\circ \cdots \circ \psi_i : F_i\otimes \pi^*K_C^{-m} \to F_i\]
    On the other hand, the map $\nu_2: \pi^*K_C^{-m}\to A_Y(m)$ is defined by multiplication by the tautological section $\lambda\in H^0(Y, \pi^*K_C^m)$. 
    %\[ F\otimes_{\O_Y} \pi^*K_C^{-m} \to F\otimes_{\O_Y} A_Y \to F, \quad s\otimes z \mapsto  \lambda(z) s \otimes 1 \mapsto \lambda(z)s \]
    Its induced action on each component $F_i$ maps a section $f\otimes z\in F_i\otimes_{\O_Y}\pi^*K_C^{-m}$ to $f\cdot \lambda(z)\in F_i$, that is, 
    \[\operatorname{Id}_{F_i}\otimes _{\O_Y}\lambda: F_i\otimes \pi^*K_C^{-m} \to F_i.\]
    Therefore, in order for $\m F$ to annihilates the ideal $I$, each component $F_i$ must satisfy
    \[\psi_{i-1}\circ \cdots \circ \psi_i = \operatorname{Id}_{F_i}\otimes_{\O_Y} \lambda  ,\quad \textrm{for } i\in \zm\]
    as desired.  
\end{proof}

\end{document}
